% ============================================================
% Section 28: 振子强度求和规则 (Thomas–Reiche–Kuhn Sum Rule)
% ============================================================
\section{量子力学：振子强度求和规则——从谐振子到普遍束缚系统}

\begin{modelbox}{题目}
质量为 $M$ 的粒子在一维势场中运动，束缚态能级为 $E_n$，本征态为 $|u_n\rangle$。
定义 \textbf{振子强度} (oscillator strength)：
\[
  f_n = \frac{2M}{\hbar^2}(E_n - E_0) \bigl|\langle u_n|X|u_0\rangle\bigr|^2.
\]
\textbf{要求}：
\begin{enumerate}
  \item 对一维谐振子，计算 $\displaystyle\sum_n f_n$。
  \item 对一般一维束缚系统，计算 $\displaystyle\sum_n f_n$。
\end{enumerate}
\end{modelbox}

\vspace{0.3em}
\begin{insightbox}{结论预告}
两者结果均为 $\displaystyle\sum_n f_n = 1$。
谐振子是"恰好"贡献全部；一般系统则通过对易关系 $[X,P]=i\hbar$ 和完备性给出 \textbf{普适求和规则}，与势场形式无关。
\end{insightbox}


\subsection{一维谐振子}

\subsubsection{能级与矩阵元}

一维谐振子哈密顿量：
\[
  H = \frac{P^2}{2M} + \frac{1}{2}M\omega^2 X^2, \qquad
  E_n = \hbar\omega\Bigl(n+\frac{1}{2}\Bigr).
\]

位置算符用产生湮灭算符表示：
\[
  X = \sqrt{\frac{\hbar}{2M\omega}}\,(a + a^\dagger).
\]

\begin{derivationbox}{只有 $n=1$ 有贡献}
\[
  a|u_0\rangle = 0, \qquad a^\dagger|u_0\rangle = |u_1\rangle.
\]
因此：
\[
  \langle u_n|X|u_0\rangle
  = \sqrt{\frac{\hbar}{2M\omega}}\,
    \langle u_n|(a+a^\dagger)|u_0\rangle
  = \sqrt{\frac{\hbar}{2M\omega}}\,\delta_{n,1}.
\]
\textbf{关键观察}：$X$ 只有相邻能级间的非零矩阵元，所以 $|u_0\rangle$ 只能跃迁到 $|u_1\rangle$。
\end{derivationbox}

\subsubsection{计算 $f_n$}

对 $n=1$：
\begin{align*}
  E_1 - E_0 &= \hbar\omega, \\
  \bigl|\langle u_1|X|u_0\rangle\bigr|^2 &= \frac{\hbar}{2M\omega}.
\end{align*}

代入定义：
\[
  f_1 = \frac{2M}{\hbar^2} \cdot \hbar\omega \cdot \frac{\hbar}{2M\omega} = \boxed{1}.
\]

对 $n\neq 1$，矩阵元为零，$f_n = 0$。

\subsubsection{求和结果}

\begin{formulabox}{谐振子求和规则}
\[
  \sum_{n=0}^{\infty} f_n = f_1 = \boxed{1}.
\]
\end{formulabox}


\subsection{一般一维束缚系统}

\subsubsection{核心对易关系}

\begin{tipbox}{$[X,H]$ 给出 $P$ 的矩阵元}
\[
  [X,H] = \Bigl[X, \frac{P^2}{2M} + V(X)\Bigr]
       = \frac{1}{2M}[X,P^2]
       = \frac{i\hbar}{M}\,P.
\]
因此：
\[
  \boxed{P = \frac{M}{i\hbar}[X,H]}.
\]
\end{tipbox}

\subsubsection{$P$ 的跃迁矩阵元}

取 $|u_0\rangle$ 与 $|u_n\rangle$ 之间的矩阵元：
\begin{align*}
  \langle u_n|P|u_0\rangle
  &= \frac{M}{i\hbar}\langle u_n|[X,H]|u_0\rangle \\
  &= \frac{M}{i\hbar}(E_0 - E_n)\langle u_n|X|u_0\rangle. \tag{1}
\end{align*}

\subsubsection{利用完备性求和}

\begin{thinkbox}{为什么想到 $[X,P]=i\hbar$？}
$f_n$ 中含有 $(E_n-E_0)|\langle u_n|X|u_0\rangle|^2$。要从对易关系得到这个组合，最自然的思路是：
\begin{itemize}
  \item 用 $[X,H]$ 把 $P$ 与 $X$ 联系起来（产生能量差因子 $E_n-E_0$）
  \item 再用 $[X,P]=i\hbar$ 给出闭合条件
\end{itemize}
这是求和规则 (sum rule) 的标准推导套路。
\end{thinkbox}

考虑 $[X,P]=i\hbar$ 在基态的期望值，插入完备性关系 $\sum_n |u_n\rangle\langle u_n| = I$：
\begin{align*}
  \langle u_0|[X,P]|u_0\rangle
  &= \sum_n \Bigl(
       \langle u_0|X|u_n\rangle\langle u_n|P|u_0\rangle
       - \langle u_0|P|u_n\rangle\langle u_n|X|u_0\rangle
     \Bigr) \\
  &= i\hbar. \tag{2}
\end{align*}

利用式 (1) 及其厄米共轭 $\langle u_0|P|u_n\rangle = \langle u_n|P|u_0\rangle^* = -\frac{M}{i\hbar}(E_0-E_n)\langle u_0|X|u_n\rangle$，代入式 (2)：
\begin{align*}
  i\hbar
  &= \sum_n \frac{M}{i\hbar}(E_0-E_n)\bigl|\langle u_n|X|u_0\rangle\bigr|^2
   - \sum_n \Bigl(-\frac{M}{i\hbar}(E_0-E_n)\Bigr)\bigl|\langle u_n|X|u_0\rangle\bigr|^2 \\
  &= \frac{2M}{i\hbar}\sum_n (E_0-E_n)\bigl|\langle u_n|X|u_0\rangle\bigr|^2.
\end{align*}

两边乘以 $\frac{i\hbar}{2M}$：
\[
  \sum_n (E_n-E_0)\bigl|\langle u_n|X|u_0\rangle\bigr|^2 = \frac{\hbar^2}{2M}. \tag{3}
\]

\subsubsection{得到振子强度求和}

将式 (3) 代入 $f_n$ 的定义：
\begin{align*}
  \sum_n f_n
  &= \sum_n \frac{2M}{\hbar^2}(E_n-E_0)\bigl|\langle u_n|X|u_0\rangle\bigr|^2 \\
  &= \frac{2M}{\hbar^2} \cdot \frac{\hbar^2}{2M} \\
  &= \boxed{1}.
\end{align*}

\begin{formulabox}{Thomas–Reiche–Kuhn (TRK) 求和规则}
\[
  \boxed{\sum_n f_n = \frac{2M}{\hbar^2}\sum_n (E_n-E_0)\bigl|\langle u_n|X|u_0\rangle\bigr|^2 = 1}
\]
对\textbf{任意}一维束缚势场 $V(x)$ 均成立，与势的具体形式无关。
\end{formulabox}


\subsection{物理意义与考场策略}

\subsubsection{物理意义}

\begin{insightbox}{为什么必须是 1？}
\textbf{Thomas–Reiche–Kuhn 求和规则} 是量子力学完备性（完备基）和基本对易关系 $[X,P]=i\hbar$ 的直接推论。

它表明：从基态到所有激发态的"振子强度"总和，恰好等于一个经典谐振子的强度。这不是巧合——它反映了量子系统必须满足与经典极限的对应关系（$\hbar\to 0$ 时量子结果退化为经典结果）。

\textbf{多维推广}：对三维各向同性系统，$\sum_n f_n^{(x)} = \sum_n f_n^{(y)} = \sum_n f_n^{(z)} = 1$，因此 $x,y,z$ 方向总和为 3。
\end{insightbox}

\subsubsection{常见错误}

\begin{errorbox}{易错点}
\begin{enumerate}
  \item \textbf{漏掉 $n=0$ 项}：虽然 $E_0-E_0=0$，形式上 $f_0=0$，但书写求和时应明确是否包含 $n=0$。
  \item \textbf{符号错误}：$\langle u_n|P|u_0\rangle = \frac{M}{i\hbar}(E_0-E_n)\langle u_n|X|u_0\rangle$ 中的 $(E_0-E_n)$ 不是 $(E_n-E_0)$。
  \item \textbf{忘记厄米共轭}：$\langle u_0|P|u_n\rangle$ 要用 $\langle u_n|P|u_0\rangle^*$，带来额外负号。
  \item \textbf{误以为需要知道 $V(x)$}：TRK 规则是\textbf{普适的}，不需要解出具体波函数。
\end{enumerate}
\end{errorbox}

\subsubsection{考场速记}

\begin{cheatbox}{TRK 求和规则：考场三步法}
\textbf{Step 1}——写 $[X,H]$：
\[
  [X,H] = \frac{i\hbar}{M}P \quad\Rightarrow\quad P = \frac{M}{i\hbar}[X,H]
\]
\textbf{Step 2}——求 $P$ 矩阵元（出能量差）：
\[
  \langle u_n|P|u_0\rangle = \frac{M}{i\hbar}(E_0-E_n)\langle u_n|X|u_0\rangle
\]
\textbf{Step 3}——用 $[X,P]=i\hbar$ 闭合：
\[
  \langle u_0|[X,P]|u_0\rangle = i\hbar \xrightarrow{\text{插入完备基}} \sum_n (E_n-E_0)|\langle u_n|X|u_0\rangle|^2 = \frac{\hbar^2}{2M}
\]
\textbf{结论}：$\displaystyle\sum_n f_n = 1$。
\end{cheatbox}
